\pdfoutput=1
\documentclass[11pt]{article}
\usepackage[T1]{fontenc}
\usepackage{lmodern}
\usepackage[a4paper,margin=26mm]{geometry}
\usepackage{amsmath,amssymb,amsthm,booktabs,hyperref}
\usepackage{tikz,needspace}
\hypersetup{colorlinks=true,urlcolor=blue,citecolor=blue,
  pdftitle={The ten minimum universal sets of MacMahon cubes},pdfauthor={Jason Davies}}
\newtheorem{theorem}{Theorem}
\newtheorem{lemma}[theorem]{Lemma}
\newtheorem{corollary}[theorem]{Corollary}
\emergencystretch=2em
\newcommand{\lean}[1]{\nolinkurl{#1}}
\makeatletter\g@addto@macro\UrlBreaks{\do\_}\makeatother
\title{The ten minimum universal sets of MacMahon cubes}
\author{Jason Davies\thanks{\href{mailto:jason@jasondavies.com}{\nolinkurl{jason@jasondavies.com}}}}
\date{November 1, 2026}
\begin{document}
\maketitle
\begin{abstract}
A collection of six-coloured MacMahon cubes is universal if eight of its
cubes can be selected to build a $2\times2\times2$ model of any target
variety, with matching exterior colours and no internal matching requirement.
We prove that exactly ten labelled collections attain the known minimum
size of twelve, even when repeated varieties are allowed, resolving the
conjectures of Johnson and Rold\'an. Usable-cube inequalities force a
symmetric multiplicity matrix with row and column sums two; corner coverage
then forces two disjoint triangles in Conway's tableau. The written proof
is elementary and requires no exhaustive search. The classification and its
discrete geometric interpretation are formalised in Lean~4. Exhaustive
computation provides a separate check and additional target-count
distributions.
\end{abstract}

\noindent\textbf{AI transparency statement.} The classification proof,
computational investigation, Lean formalisation and manuscript were produced
by AI systems under the author's direction. The systems and their
contributions are described at the end of the paper.

\section{Introduction}
MacMahon's cube puzzles use the thirty varieties obtained by colouring the
six faces with six distinct colours, identifying proper rotations but not
reflections \cite[pp.~42--46]{MacMahon}. We study the exterior-matching
puzzle known as \emph{Eight Blocks to Madness}: select eight cubes and
assemble a $2\times2\times2$ model of a target variety. Only the exterior
colours must match; touching interior faces are unrestricted.
A collection is \emph{universal} if it can build every target.
Collections may contain repeated varieties; a \emph{set} contains distinct
varieties. Face-colour labels remain fixed throughout.

Haraguchi proved that twelve is the minimum size of a universal collection,
allowing repeats, and exhibited one attaining it
\cite[Lemma~2 and Theorem~4]{H}. Johnson and Rold\'an found nine further
sets, bringing the known total to ten \cite[Theorem~3]{JR}. Their Conjecture~4 asserts that these are all the
minimum sets; Conjecture~8 describes them as two triangles in Conway's
tableau. This tableau labels varieties by ordered pairs $ij$ of distinct
labels in $L=\{A,B,C,D,E,F\}$, as specified in Section~\ref{sec:tableau}.

\begin{theorem}[Minimum universal collections]\label{thm:main}
The minimum size of a universal collection is twelve. Exactly ten labelled
collections attain that size, even when repeated varieties are permitted.
Each unordered partition $L=P\sqcup Q$ into triples gives one such
collection: one copy of each variety $ij$ with $i\ne j$ and both labels
in the same part. These are all the minimum collections; in particular, none has repeats.
\end{theorem}

Our contribution is an elementary proof of this classification, its extension
to repeated varieties, and a Lean formalisation of the classification
theorems. The written proof first uses usable-cube inequalities to force
symmetric multiplicities and row sums two. The resulting support graphs are
either two triangles or bipartite; corner coverage excludes the latter.
An explicit construction proves sufficiency. The formal proof uses a different
route for these last two steps, described in Section~\ref{sec:formal}.
We also answer the computational questions in \cite[Section~6]{JR}: for
sets of nine, ten and eleven distinct varieties the maximum target counts
are respectively 9, 20 and 22, and thirteen is the smallest set size
that permits exactly 29 targets.

\section{Conventions and tableau facts}\label{sec:tableau}
Each cube has six different face colours, labelled $1,\ldots,6$.
There are $6!/24=30$ rotational varieties. The ten minimum sets are counted
separately even though they form one colour-permutation orbit.

Our notation $ij$ corresponds to the conventional upper/lower-case cube
name, such as $Ab$ \cite[Section~3]{B2017}. The tableau labels in $L$ are
distinct from the face-colour labels $1,\ldots,6$. At each vertex, the three incident face colours define an
\emph{oriented corner type}: an ordered triple modulo cyclic permutation,
read clockwise when viewed from outside the cube.
Reversing the order gives a different type.
Each cube has eight corners, and there are $2\binom63=40$ types in total.

Table~\ref{tab:cubes} lists face colours in the order $+x,-x,+y,-y,+z,-z$. The frame $(+x,+y,+z)$ is
left-handed, so the following rule gives clockwise corner triples.
Choose one face from each opposite pair and list their colours in axis
order $x,y,z$; reverse the order when an odd number of chosen faces have
negative signs. Reduce each triple modulo cyclic permutation. This recovers
the eight corner types of every entry directly from the table.

\begin{table}[htbp]
\centering
\begin{tabular}{lrrrrrr}\toprule
 & $a$ & $b$ & $c$ & $d$ & $e$ & $f$ \\\midrule
$A$ & $-$ & \texttt{152463} & \texttt{162354} & \texttt{142635} & \texttt{132546} & \texttt{123456} \\
$B$ & \texttt{152436} & $-$ & \texttt{123564} & \texttt{162543} & \texttt{142365} & \texttt{132654} \\
$C$ & \texttt{162345} & \texttt{123546} & $-$ & \texttt{132456} & \texttt{152634} & \texttt{142563} \\
$D$ & \texttt{142653} & \texttt{162534} & \texttt{132465} & $-$ & \texttt{123645} & \texttt{152346} \\
$E$ & \texttt{132564} & \texttt{142356} & \texttt{152643} & \texttt{123654} & $-$ & \texttt{162435} \\
$F$ & \texttt{123465} & \texttt{132645} & \texttt{142536} & \texttt{152364} & \texttt{162453} & $-$ \\
\bottomrule\end{tabular}
\caption{Face-colour representatives of Conway's tableau. Each six-digit
entry lists the colours on $+x,-x,+y,-y,+z,-z$ in a left-handed frame.
The row and column give the cube name, for example $Ab=152463$.
Their corner types agree with \cite[Table~3]{JR}.}
\label{tab:cubes}
\end{table}


\begin{lemma}[Corner--directed-triangles correspondence]
\label{lem:corner-triangles}
The forty oriented corner types are in bijection with partitions of $L$
into two triples, each equipped with a directed 3-cycle. The six cube
varieties containing a given corner are exactly the six arcs of its two
cycles. Equivalently, cube $ij$ contains a corner precisely when the corner's
two cycles contain the arc $i\to j$.
\end{lemma}
\begin{proof}
Table~\ref{tab:corners} records the correspondence explicitly. For a triple
written $XYZ$ in increasing label order, write $+$ for the cycle
$X\to Y\to Z\to X$ and $-$ for its reverse. The first and second signs
refer to the left and right triples respectively.

To verify Table~\ref{tab:corners}, compute the corners of the representatives
in Table~\ref{tab:cubes}, list the cube names containing each corner, and
read each name as an arc. For example, the $(-x,+y,+z)$ corner of $Ab$
has colours $(5,2,6)$ in axis order. Reversing gives $(5,6,2)$, or $256$
modulo cyclic permutation. The full list of cubes with corner $256$ is
\[
 Ab,\ Bc,\ Ca,\ De,\ Ef,\ Fd.
\]
These give $A\to B\to C\to A$ and $D\to E\to F\to D$, so $256$
occupies the $(+,+)$ column of the $ABC\mid DEF$ row. For every entry the
same check gives precisely the six arcs specified by its row and column.
The forty entries are distinct corner types, and the ten rows and four
columns exhaust the partitions and direction choices. This verifies the
bijection and the incidence rule. In particular, an arc $i\to j$ belongs
to eight cycle pairs: four choices of the third label in its cycle and two
directions for the complementary cycle, giving the eight corners of $ij$.
\end{proof}

\begin{table}[ht]
\centering
\begin{tabular}{lrrrr}\toprule
Partition & $(+,+)$ & $(+,-)$ & $(-,+)$ & $(-,-)$ \\\midrule
$ABC\mid DEF$ & 256 & 143 & 134 & 265 \\
$ABD\mid CEF$ & 465 & 132 & 123 & 456 \\
$ABE\mid CDF$ & 345 & 126 & 162 & 354 \\
$ABF\mid CDE$ & 235 & 164 & 146 & 253 \\
$ACD\mid BEF$ & 364 & 125 & 152 & 346 \\
$ACE\mid BDF$ & 356 & 142 & 124 & 365 \\
$ACF\mid BDE$ & 153 & 246 & 264 & 135 \\
$ADE\mid BCF$ & 165 & 243 & 234 & 156 \\
$ADF\mid BCE$ & 136 & 254 & 245 & 163 \\
$AEF\mid BCD$ & 263 & 145 & 154 & 236 \\
\bottomrule\end{tabular}
\caption{Corner types by partition and cycle directions. On each displayed
triple $XYZ$, $+$ means $X\to Y\to Z\to X$ and $-$ means its reverse;
the two signs specify the left and right triples. Corner numbers use the
fixed face-colour labels of \cite{JR}, modulo cyclic permutation.}
\label{tab:corners}
\end{table}

The standard corner-sharing rule \cite[Lemma~2.7]{B2008}, in tableau form
\cite[Propositions~1--2 and Properties~1--2]{H}, follows from
Lemma~\ref{lem:corner-triangles}: two distinct varieties $ij$ and $uv$ share
no corner exactly when
$i=u$, $j=v$, or $(u,v)=(j,i)$. In every other case they share two corners.
Indeed, two distinct arcs in the same directed 3-cycle must form a directed
path; this fixes that cycle, leaving two orientations of the other one.
For two disjoint arcs, the remaining two labels can be assigned to their
respective triples in two ways, each determining both cycle directions.
Common tails, common heads, and reversed arcs cannot coexist in directed
3-cycles. The target itself shares all eight corners.

Call a cube \emph{usable} for a target when they share an oriented corner.
Each cube is usable for $30-9=21$ targets, including itself. A target can
be built exactly when its eight corners can be assigned distinct physical
cubes containing the respective corner types: an oriented-corner match
permits a proper rotation matching all three exposed faces. Internal faces
impose no further conditions.

\section{Necessity}
Let $a_{ij}$ be the number of copies of variety $ij$ in a collection of size
$K$, with $a_{ii}=0$. Write $r_i=\sum_j a_{ij}$ and $c_j=\sum_i a_{ij}$.

\Needspace{10\baselineskip}
\begin{lemma}[Row bound]
\label{lem:row-bound}
If every target has at least eight usable physical cubes in the collection,
then
\[
 r_i+c_j-2a_{ij}+a_{ji}\le K-8\qquad(i\ne j),
 \qquad 3r_i\le 4K-40,\quad 3c_i\le4K-40.
\]
\end{lemma}
\begin{proof}
For target $ij$, the unusable cubes are those in its row apart from $ij$,
those in its column apart from $ij$, and the copies of $ji$. These classes
are disjoint and total $r_i+c_j-2a_{ij}+a_{ji}$. At least eight cubes must
remain usable, proving the first inequality. Summing over the five
$j\ne i$ gives
\[
5r_i+(K-c_i)-2r_i+c_i=3r_i+K\le5(K-8).
\]
The column bound follows by the analogous sum.
\end{proof}

\begin{lemma}[Necessary support]
\label{lem:necessary-support}
A universal collection of size twelve consists of one copy of every directed
edge within each of two disjoint triples of tableau labels.
\end{lemma}
\begin{proof}
Universality supplies eight usable cubes for each target, so
Lemma~\ref{lem:row-bound} applies and gives $r_i\le\lfloor8/3\rfloor=2$.
As the six row sums total twelve, all equal two. The same holds for the columns. If $a_{ij}=0$, the
first inequality now gives $4+a_{ji}\le4$, so $a_{ji}=0$. Thus support is
symmetric.

Consider the undirected support graph on the six tableau labels $L$: an
edge $\{i,j\}$ is present when $a_{ij}>0$. Every vertex has degree one
or two. If $i$ has only one neighbour $j$, its row sum gives $a_{ij}=2$.
By support symmetry, column $i$ has no other nonzero entry, so $a_{ji}=2$.
These entries exhaust both row sums: $i$ and $j$ form an isolated pair,
with two copies in each direction. We call this a \emph{double-copy pair}.
All other vertices have degree two, forming simple cycles with one copy
in each direction on every edge.

On six vertices the only possibilities are two triangles, a 6-cycle,
a 4-cycle and a double-copy pair, or three double-copy pairs. All except
two triangles are bipartite with parts of size three. Choose these parts
as the two triples defining a corner type. Every directed edge of that
corner lies within a part, whereas every selected variety joins the parts.
Thus no cube in the collection has that corner type. A target having
that corner cannot be built, contradicting universality. Only two triangles
remain, and their arc multiplicities are all one.
\end{proof}

For distinct varieties, the degree argument leaves only a 6-cycle or two
triangles; Figure~\ref{fig:supports} illustrates why the 6-cycle fails.
\begin{figure}[htbp]
\centering
\input{support-graphs.tex}
\caption{The two possible degree-two supports for twelve distinct varieties.
Each undirected edge represents both cube varieties $ij$ and $ji$, hence
six edges represent twelve cubes. Circles mark $\{A,B,C\}$ and squares
mark $\{D,E,F\}$; these are tableau labels, not face colours. The left
support misses all four corners for this partition (Table~\ref{tab:corners});
the right support is universal. Repeated varieties allow the additional
support types treated in the proof.}
\label{fig:supports}
\end{figure}

\Needspace{10\baselineskip}
\section{Sufficiency and classification}
\begin{lemma}[Two-triangle construction]
\label{lem:sufficiency}
For any partition $L=P\sqcup Q$ into triples, the twelve directed edges
within $P$ and $Q$ form a universal set.
\end{lemma}
\begin{proof}
Fix target $ij$ and interchange $P,Q$ if necessary so that $i\in P$.
By Lemma~\ref{lem:corner-triangles}, its corners are indexed by a third
label $k\notin\{i,j\}$ and either direction on the complementary triple;
the cycle containing $ij$ is $i\to j\to k\to i$.

If $j\in P$, write $P=\{i,j,p\}$. For each $k\in Q$, the
complementary cycle contains one arc between the two labels of $Q\setminus
\{k\}$. Assign that cube to the corner. As $k$ and the cycle direction vary,
these are the six distinct directed edges within $Q$. The remaining two
corners have $k=p$; fill them with $jp$ and $pi$, both of which belong to
both corner types. This gives eight distinct cubes.

If $j\in Q$, write $P\setminus\{i\}=\{u,v\}$ and
$Q\setminus\{j\}=\{x,y\}$. For the four corners with $k\in\{u,v\}$,
the selected varieties containing each such corner are $ki$ and one of
$xy,yx$. Their incidence
is described by a $K_{2,2}$ with vertex classes $\{ui,vi\}$ and
$\{xy,yx\}$. Here the four corners are the edges, and the four candidate
cube varieties are the vertices: each corner joins its two candidates.
Orient the resulting 4-cycle cyclically and assign each edge to its head.
Every vertex then receives exactly one edge, so each corner receives an
incident candidate and all four assigned cubes are distinct.
For the other four corners, the analogous four varieties are $jx,jy,uv,vu$.
The two groups of four varieties are disjoint, completing the construction.
\end{proof}

\begin{proof}[Proof of Theorem~\ref{thm:main}]
Each physical cube is usable for 21 targets, so universality requires
$21K\ge30\cdot8=240$, hence $K\ge12$. This lower bound is due to
Haraguchi \cite[Lemma~2]{H}. The preceding lemmas give necessity and sufficiency at
twelve. There are $\frac12\binom63=10$ unordered partitions into triples.
\end{proof}

\begin{corollary}[Twelve-cube criterion]
\label{cor:twelve-criterion}
A collection of twelve physical cubes, with repeated varieties allowed,
is universal if and only if both of the following hold:
\begin{enumerate}
\item Every target has at least eight usable cubes in the collection,
counting physical copies.
\item Every oriented corner type occurs on at least one cube in the
collection.
\end{enumerate}
\end{corollary}
\begin{proof}
Universality supplies eight usable cubes for each target. Every oriented
corner type belongs to some target, so it must occur in the collection.

Conversely, condition~(1) gives the inequalities of
Lemma~\ref{lem:row-bound}. At size twelve these force every row and column
sum to equal two and the support to be symmetric, exactly as in the proof
of Lemma~\ref{lem:necessary-support}. The same degree argument leaves the
four support types listed there. Each bipartite type omits all four corners
associated with its balanced bipartition, contradicting condition~(2).
Only the two-triangle type remains, with one copy of each arc, and
Lemma~\ref{lem:sufficiency} proves it universal.
\end{proof}

The six-cycle collection in Figure~\ref{fig:supports} shows why condition~(2)
is needed. Its row and column sums are two and $a_{ij}=a_{ji}$, so target
$ij$ has $12-(r_i+c_j-2a_{ij}+a_{ji})=8+a_{ij}\ge8$ usable cubes.
Thus condition~(1) holds, but the four corners associated with its balanced
bipartition are absent.

\section{Computational checks and additional results}
An exhaustive enumeration of all $\binom{30}{12}=86,493,225$ subsets returns
exactly the ten sets above. The implementation also reproduces the published
distribution in \cite[Table~2]{JR} for all
$\binom{30}{8}=5,852,925$ eight-cube sets.
Two algorithms implement Haraguchi's criteria: bipartite matching
\cite[Proposition~3]{H} and the tree-component test
\cite[Proposition~4]{H}. The two feasibility tests share the subset
enumerator and tableau. The latter test uses a compatibility multigraph on all
eight target corners, including isolated vertices. Each usable non-target
physical cube is an edge joining its two compatible corners; parallel
edges are allowed. Here cubes are edges and corners are vertices, whereas
the necessity proof uses a support graph on six tableau labels.
A tree component, including an isolated vertex, leaves one corner unfilled
by its edges. A component with a cycle can fill all its corners; two
parallel edges already form a cycle. Since a copy of the target can fill
any corner, the target is buildable exactly when the number of tree
components is at most the number of target copies. The subset enumeration
precomputes feasibility for all $2^{20}$ subsets of the twenty usable
non-target varieties for each target, and uses no symmetry reduction.

Separate geometric checks generate the 24 proper rotations, 30 varieties
and 40 corner types from face-colour permutations. They verify the tableau,
the directed-triangle correspondence and physical placements for all
300 minimum-set--target pairs. A separate enumeration of the 70 simple
2-regular graphs gives 60 six-cycles, each building 18 targets, and the ten
universal triangle-pair sets.

The single colour-permutation orbit observed in \cite[Section~6]{JR} also
follows from the classification. Colour permutations induce the full
symmetric group $S_6$ on the six tableau rows \cite[Lemma~3.1 and
Corollary~3.3]{B2017}. Since colour permutations commute with reflection,
which sends $ij$ to $ji$, rows and columns are permuted identically.
Thus the action on entries is $ij\mapsto\pi(i)\pi(j)$, and $S_6$ acts
transitively on unordered partitions into two triples. The stabiliser of a partition has order $2(3!)^2=72$, since
it may permute each triple and interchange the two triples.

For sets of distinct varieties the exhaustive target-count ranges are:
\[
\begin{array}{c|rrrrrr}
\text{set size}&8&9&10&11&12&13\\\hline
\text{minimum targets}&0&0&1&2&8&10\\
\text{maximum targets}&5&9&20&22&30&30
\end{array}
\]
The lower bound of one target for ten distinct varieties was proved in
\cite[Lemma~4.1]{B2017}; the enumeration confirms that it is sharp.
No set of size at most twelve builds exactly 29 targets. At size thirteen,
827,250 sets build exactly 29 targets and 53,100 are universal.
These enumeration results answer the questions in \cite[Section~6]{JR}
about the maxima for sizes nine, ten and eleven and the possibility of
building exactly 29 targets. The earlier distributions for sizes nine
through twelve used random samples of 20,000 sets each; our ranges are
exhaustive. These target-count bounds rely on the complete enumeration.
Full histograms, example sets, source code and checkers are in the
accompanying repository.\footnote{\url{https://github.com/jasondavies/macmahon-min-universal-sets}}

\section{Formal verification and reproducibility}\label{sec:formal}
The accompanying repository includes a formalisation using Lean~4.35.0
\cite{Lean4} and mathlib v4.35.0 \cite{mathlib}. The main results are
\lean{MacMahon.minimum_size}, \lean{MacMahon.physical_exactly_ten},
and \lean{MacMahon.physical_iff_two_triangles}.

In Lean, an inventory records a natural-number multiplicity for each tableau entry.
The stored face-colour representatives are those in Table~\ref{tab:cubes},
with the same face order and handedness.
A construction assigns eight distinct physical copies to the eight positions,
with proper rotations matching the exposed faces. Lean proves that the
rotation table consists exactly of determinant-positive signed-axis
permutations, that the thirty representatives and their rotations give all
720 face-colour assignments without duplication, and that corner matching
is equivalent to exposed-face matching.

The formalisation proves the classification theorems by a different route
from the written support-graph and construction arguments. It shares the
usable-cube inequalities and symmetry argument, allowing arbitrary
multiplicities. For necessity, a label permutation puts the first row into one of
two forms: a doubled edge or two distinct neighbours. Arithmetic proofs
exclude the doubled edge and force the neighbours to close a triangle;
the degree equations then determine the complementary triangle. For
sufficiency, Lean checks explicit witnesses for all 300 minimum-set--target
pairs instead of formalising the general construction of
Lemma~\ref{lem:sufficiency}. These finite checks use kernel reduction. The proof library has no unfinished
proofs or added axioms: audits permit only propositional extensionality,
classical choice and quotient axioms. The generated tables and construction
witnesses are certified by kernel-checked Lean proofs; the data generator
itself is not trusted. The target-count enumerations
are outside the formalisation, and Corollary~\ref{cor:twelve-criterion}
has a written proof but no separate Lean statement.

The development is packaged for the comparator of the Palomar registry.
Ten independent statements in
\texttt{Challenge.lean}, which imports only mathlib, are separated from
completed proofs in \texttt{Solution.lean}. The comparator checks statement
equality, permitted axioms and proof acceptance by Lean's kernel and the
independent kernels NanoDa and con-ron. Reproduction commands are in the
accompanying repository; local verification does not constitute registry
acceptance.

\section*{Declaration of AI-generated research and review}
GPT-6-Astra (OpenAI),
used through Codex, developed the classification proof and computational
investigation, wrote the Lean formalisation, and wrote and revised the
manuscript. Claude Opus~5.5 (Anthropic), used through Claude Code, contributed to
paper review.

The formal claims are checked by Lean and the additional kernels described
in Section~\ref{sec:formal}; the target-count computations have separate
reproducible checks. These automated checks are not independent human peer
review. The literature search was limited to the cited sources and related
material located during development; publication priority has not been
established by a comprehensive review.

The human author initiated and directed the research and accepts
responsibility for the content of this paper and the accuracy of this
disclosure.

\begin{thebibliography}{9}
\raggedright
\bibitem{MacMahon} P.~A.~MacMahon,
\emph{New Mathematical Pastimes}, Cambridge University Press (1921).
\href{https://archive.org/details/newmathematicalp00macmuoft}{Digitised edition}.
\bibitem{JR} I.~Johnson and \'E.~Rold\'an,
\emph{Solution Numbers for Eight Blocks to Madness Puzzle},
\href{https://arxiv.org/abs/2407.13208v1}{arXiv:2407.13208v1} (2024).
\bibitem{H} K.~Haraguchi,
\emph{On a generalization of ``Eight Blocks to Madness'' puzzle},
Discrete Mathematics \textbf{339}(4) (2016), 1400--1409.
\href{https://doi.org/10.1016/j.disc.2015.12.014}{doi:\nolinkurl{10.1016/j.disc.2015.12.014}}.
Preprint \href{https://arxiv.org/abs/1408.3696v2}{arXiv:1408.3696v2};
numbered results cited here follow this version.
\bibitem{B2017} E.~Berkove, D.~Cervantes Nava, D.~Condon and R.~Katz,
\emph{Automorphisms of $S_6$ and the color cubes puzzle},
Australasian Journal of Combinatorics \textbf{68}(1) (2017), 71--93.
\href{https://ajc.maths.uq.edu.au/pdf/68/ajc_v68_p071.pdf}{Publisher PDF}.
\href{https://arxiv.org/abs/1503.07184}{arXiv:1503.07184}.
\bibitem{B2008} E.~Berkove, J.~Van Sickle, B.~Hummon and J.~Kogut,
\emph{An analysis of the (colored cubes)$^3$ puzzle},
Discrete Mathematics \textbf{308}(7) (2008), 1033--1045.
\href{https://doi.org/10.1016/j.disc.2007.03.058}{doi:\nolinkurl{10.1016/j.disc.2007.03.058}}.
\bibitem{Lean4} L.~de Moura and S.~Ullrich,
\emph{The Lean 4 Theorem Prover and Programming Language},
in Automated Deduction---CADE~28, LNCS \textbf{12699}, Springer (2021),
625--635.
\href{https://doi.org/10.1007/978-3-030-79876-5_37}{doi:\nolinkurl{10.1007/978-3-030-79876-5_37}}.
\bibitem{mathlib} The mathlib Community,
\emph{The Lean Mathematical Library},
in Proceedings of the 9th ACM SIGPLAN International Conference on
Certified Programs and Proofs (CPP~2020), ACM (2020), 367--381.
\href{https://doi.org/10.1145/3372885.3373824}{doi:\nolinkurl{10.1145/3372885.3373824}}.
\href{https://arxiv.org/abs/1910.09336}{arXiv:1910.09336}.
\end{thebibliography}
\end{document}
